\documentclass[12pt]{article}
\usepackage[utf8]{inputenc}
\usepackage[T1]{fontenc}
\usepackage{hyperref}
\title{Sulka Language Resources for K-12 Curriculum Development}
\author{Compiled by Claude Sonnet 5.5 (claude-sonnet-5-5) for Andr\'as Kornai}
\date{}
\begin{document}
\maketitle

\section{Introduction}
Sulka (ISO 639-3: sua; Glottocode: sulk1246) is a language isolate spoken in the eastern Pomio District, East New Britain Province, Papua New Guinea. Published estimates of speaker numbers range from about 2,500 to 3,500. This report rests on a single web search and is deliberately short.

\section{Primary Educational Resources}
No primers, textbooks or school programmes for Sulka were found in the search results.

\section{Linguistic and Reference Materials}
\begin{itemize}
\item Reesink, G. (2005), ``Sulka of East New Britain: A Mixture of Oceanic and Papuan Traits'', \emph{Oceanic Linguistics} 44(1): 145--193 (Project MUSE): \url{https://muse.jhu.edu/article/184769/}
\item Grambank entry for Sulka: \url{https://grambank.clld.org/languages/sulk1246}
\item Grambank source record for Reesink's Sulka description: \url{https://grambank.clld.org/sources/sReesinkSulka}
\item Glottolog reference record for Clifton (1996), SIL, \emph{Two non-Austronesian grammars from the islands} (Kuot and Sulka): \url{https://glottolog.org/resource/reference/id/142755}
\item Wikipedia overview of the Sulka language: \url{https://en.wikipedia.org/wiki/Sulka_language}
\end{itemize}
Schneider (1962 [1942]), \emph{Grammatik der Sulka-Sprache (Neubritannien)}, is listed by Grambank. No Sulka dictionary was found.

\section{Religious and Community Resources}
None found.

\section{Digital and Audio Resources}
\begin{itemize}
\item Endangered Languages Project entry for Sulka: \url{https://endangeredlanguages.com/node/110541}
\end{itemize}
No audio, video or software resources were found.

\section{Curriculum Development Context}
Only academic grammars exist in the search results, and no classroom materials or dictionary. A curriculum would need new documentation work with the Sulka community.
\end{document}